% ============================================================
% Submission Checklist — Project 3 (PDE IBVP), end of project (Tutorial 13)
% Compile: tectonic -X compile submission_checklist_ibvp.tex
% ============================================================
\documentclass[11pt]{article}
\usepackage[a4paper,margin=2.3cm]{geometry}
\usepackage{amsmath,amssymb}
\usepackage{booktabs}
\usepackage{enumitem}
\setlength{\parindent}{0pt}
\setlength{\parskip}{0.5em}

\begin{document}

\begin{center}
  {\LARGE\bfseries Submission Checklist}\\[0.4em]
  {\large Project 3 (PDE IBVP) --- End of Project (Tutorial 13)}
\end{center}
\vspace{0.2em}
\hrule
\vspace{0.6em}

Go through this checklist in the last tutorial \emph{before} submitting. The
instructor runs the same checks on a fresh machine.

% ============================================================
\subsection*{The three-part submission package}
% ============================================================
\begin{itemize}[leftmargin=2.4em, itemsep=1pt]
  \item[$\square$] \textbf{Team report} --- PDF, $\geq$ 8 pages, including the
  team AI Transparency Log and one labelled ICS per member as appendices
  \item[$\square$] \textbf{Code} --- zip archive with a README
  \item[$\square$] \textbf{Slides} --- PDF, one set for the scheduled presentation wave
\end{itemize}

% ============================================================
\subsection*{Report checks}
% ============================================================
\begin{itemize}[leftmargin=2.4em, itemsep=1pt]
  \item[$\square$] Title page with team number, member names, and student IDs
  \item[$\square$] Structure is sensible, e.g.\ following the template:
    Background $\to$ Model $\to$ Schemes $\to$ Stability $\to$ Fidelity $\to$
    Conclusions
  \item[$\square$] Complete IBVP stated: PDE, domain, coefficients, units or
    nondimensionalisation, initial data, boundary data, compatibility
    assumptions, and the quantity of interest
  \item[$\square$] Two schemes derived from the PDE, including
    boundary enforcement and start-up for multistep methods; a citation may
    support a derivation but does not replace it; any library call is
    separately justified and shown to do what you claim on this problem
  \item[$\square$] Every figure has a caption saying what is shown AND why it matters
  \item[$\square$] Every claim has support (figure, table, or code output)
  \item[$\square$] A stability argument matched to your chosen pathway (von
    Neumann, matrix/energy, or conservation/entropy) stated with its
    assumptions; an experiment capable of \emph{contradicting} it was run
    and the outcome reported (a disagreement is investigated ---
    implementation, measurement, or assumption --- not asserted away)
  \item[$\square$] \emph{If your scheme is Crank--Nicolson:}
    sign-alternation ($G \to -1$) correctly diagnosed as a bounded
    oscillation, \emph{not} instability. Where the chosen investigation
    suppresses the early transient, its remedy (e.g.\ a Rannacher
    start-up with a stated tolerance) is documented --- suppression,
    not elimination (the sign alternation of the stiffest surviving
    mode remains); a team whose finding is the oscillation itself
    reports it without adding a remedy. Implicit is not automatically
    unconditionally stable: state the method, discrete operator, and
    assumptions for which that status is established (e.g.\ backward
    Euler or Crank--Nicolson for the heat equation $u_t = u_{xx}$ with
    $L_h$ symmetric negative semidefinite); only then has the scheme no
    stability boundary to cross; a conservation-law project checks the
    entropy condition instead --- the checklist item is conditional on
    the pathway, not universal
  \item[$\square$] \emph{Separated} refinements: spatial order at a fixed tiny
    $\Delta t$ (or a stated, justified error-separation protocol); temporal
    order at a fixed fine grid --- \emph{not} diagonal refinement; what is
    fixed is stated
  \item[$\square$] At least one PDE-specific fidelity diagnostic, quantified:
    decay/positivity (diffusion), mass or total variation (transport), or
    energy/phase (waves)
  \item[$\square$] One smooth reference test and one deliberately stressful
    test (nonsmooth data, long time, high frequency, or a parameter limit)
  \item[$\square$] A test capable of contradicting the stability claim was
    run and its outcome reported: show the failure case if one occurred
    (blow-up, oscillation, or a violated bound), otherwise report the
    stress test's supporting outcome with the diagnostic shown
  \item[$\square$] Your numbers are checked against your own constructed oracle
    (analytic solution, a second independent scheme, or a conserved quantity)
    --- and you state which
  \item[$\square$] Cost measure stated and reported (LO D): the two schemes
    compared at a \emph{matched} accuracy target, with the cost named (wall-clock
    time, iteration count, or operation count) and measured, not asserted
  \item[$\square$] All tables/figures state method, $h$, $\Delta t$, final
    time, norm/diagnostic, and reference
  \item[$\square$] Limitations include discretisation, modelling/domain, and
    computational limitations
  \item[$\square$] AI Transparency Log present: two or three significant claim
    audits include assumptions, risk/falsifier, independent evidence, decision,
    and limitations (or an explicit declaration that no AI was used)
  \item[$\square$] Team-owned methods and trusted library interfaces are identified
  \item[$\square$] ICS included: each member's specific contributions, traceable to git
  \item[$\square$] References cited where standard results are used
\end{itemize}

% ============================================================
\subsection*{Code checks (run on a machine with dependencies installed)}
% ============================================================
\begin{itemize}[leftmargin=2.4em, itemsep=1pt]
  \item[$\square$] \texttt{python code/run\_all.py} regenerates every figure in the report
  \item[$\square$] README gives environment setup plus one command for tests
    and one for the reported outputs
  \item[$\square$] No absolute paths that break on another machine; no caches,
    auxiliary files, or secrets included
  \item[$\square$] The model's operators, initial data, and boundary data are
    separable (not hard-coded together), so the scheme is reusable
  \item[$\square$] Multistep start-up step is implemented and correct
  \item[$\square$] Seeds fixed for any random elements
  \item[$\square$] Generated figures are reproducible (no stale \texttt{.png} from an old run)
  \item[$\square$] \texttt{.gitignore} excludes \texttt{.venv/}, \texttt{\_\_pycache\_\_/},
    build artifacts
\end{itemize}

% ============================================================
\subsection*{Slides checks}
% ============================================================
\begin{itemize}[leftmargin=2.4em, itemsep=1pt]
  \item[$\square$] One PDF set for the scheduled presentation wave
  \item[$\square$] The deck tells the full story arc (background $\to$ model $\to$ schemes $\to$ stability/fidelity $\to$ conclusions)
  \item[$\square$] Each slide has a takeaway, not just bullet points
  \item[$\square$] Presentation fits the time limit (10 min + 3 min Q\&A)
\end{itemize}

% ============================================================
\subsection*{Last sanity question}
% ============================================================
\begin{quote}
If you were the instructor, would this submission convince you?

If not, what is the weakest link --- and can you fix it in the next 30 minutes?
\end{quote}
Fix that first. A half-finished verification section is the most common reason
a project grade drops; a stability claim made outside its assumptions is the
second.

\medskip
Final reproducer/name/date: \hrulefill

\end{document}
